\documentclass[11pt, a4paper]{article}
\usepackage[utf8]{inputenc}
\usepackage[T1]{fontenc}
\usepackage{amsmath, amssymb}
\usepackage{geometry}
\usepackage[spanish]{babel}
\usepackage{fancybox}

\geometry{margin=1in}
\setlength{\parindent}{0pt}
\setlength{\parskip}{1ex plus 0.5ex minus 0.2ex}

\title{Guía de Laboratorio: Exploración de Operaciones Discretas con DSP Visualizer}
\author{IEE2103 -- Señales y Sistemas}
\date{\today}

\begin{document}
\maketitle

\section{Introducción}
Esta guía de laboratorio tiene como objetivo reforzar los conceptos teóricos de \textbf{Transformada de Fourier de Tiempo Discreto (DTFT)}, \textbf{Expansión, Diezmado, e Interpolación} mediante la herramienta interactiva \texttt{DSP Visualizer}. La aplicación les permitirá visualizar los efectos de las operaciones de cambio de tasa de muestreo directamente en los dominios del tiempo y la frecuencia ($\text{Hz}$ y $\omega$).

\begin{center}
\fbox{\parbox{0.95\textwidth}{
\textbf{Objetivos Específicos:}
\begin{itemize}
    \item Confirmar la propiedad de las imágenes espectrales generadas por la Expansión (Upsampling).
    \item Demostrar el fenómeno de Aliasing (solapamiento) causado por el Diezmado sin filtrado.
    \item Entender la función del filtro anti-aliasing y anti-imagen en el Remuestreo racional ($L/M$).
\end{itemize}
}}
\end{center}

\section{Actividad 1: La Propiedad de las Imágenes y la Interpolación ($\mathbf{L}$)}

Configure la aplicación en modo \textbf{Expansión (L)} con los siguientes parámetros iniciales:
\begin{itemize}
    \item $F_s = 200 \text{ Hz}$
    \item $f_0 = 40 \text{ Hz}$
    \item Factor $L = 3$
\end{itemize}

\subsection{Preguntas y Tareas}

\begin{enumerate}
    \item \textbf{Verificación en $\mathbf{Hz}$:} Con $L=3$, ¿cuál es la nueva frecuencia de muestreo $F_s'$?
    \item \textbf{Imágenes Espectrales:} Observe el gráfico "Espectro Original y Efecto de la Operación" (AX3). La línea roja punteada muestra el espectro de la señal expandida.
    \begin{itemize}
        \item ¿Cuántas copias del espectro original (imágenes) observa en el rango $[0, F_s']$?
        \item Indique la posición en $\text{Hz}$ de la primera imagen lateral (el primer pico positivo más allá de $f_0$). Explique cómo se relaciona esta posición con $F_s$.
    \end{itemize}
    \item \textbf{Rol del Filtro:} El proceso de interpolación (\texttt{DSP Visualizer} lo simula con \texttt{resample\_poly}) aplica un filtro ideal.
    \begin{itemize}
        \item ¿Cuál es la frecuencia de corte $\omega_c$ normalizada que debería tener este filtro para eliminar las imágenes? (Deje $L=3$).
        \item En el gráfico de Salida (DFT), ¿por qué solo se ve el pico de $40 \text{ Hz}$ y no las imágenes?
    \end{itemize}
\end{enumerate}

\section{Actividad 2: Diezmado, Aliasing y el Filtro Anti-Aliasing ($\mathbf{M}$)}

Configure la aplicación en modo \textbf{Diezmado (M)} con los siguientes parámetros iniciales:
\begin{itemize}
    \item $F_s = 200 \text{ Hz}$
    \item $f_0 = 70 \text{ Hz}$
    \item Factor $M = 3$
\end{itemize}

\subsection{Preguntas y Tareas}

\begin{enumerate}
    \item \textbf{Frecuencia de Nyquist:} Con $M=3$, calcule la nueva frecuencia de muestreo $F_s'$. ¿Cuál es la nueva frecuencia de Nyquist ($F_s'/2$)?
    \item \textbf{Observación del Aliasing:} En la sección de Diezmado, habilite la opción de simular el diezmado \textbf{SIN FILTRO}.
    \begin{itemize}
        \item ¿El pico original de $70 \text{ Hz}$ es visible en la DFT final (AX4)? ¿Qué valor de frecuencia "alias" aparece en su lugar?
        \item Explique el origen matemático de esta frecuencia de alias ($\hat{f}_0$) usando la fórmula de solapamiento: $\hat{f}_0 = |f_0 - k F_s'|$ (donde $k$ es un entero).
    \end{itemize}
    \item \textbf{Efecto del Filtro:} Vuelva a usar la función normal de Diezmado (que incluye el filtro anti-aliasing).
    \begin{itemize}
        \item ¿Qué sucede con el pico de $70 \text{ Hz}$ en el espectro de salida? ¿Por qué el filtro lo elimina?
        \item ¿Cuál es la frecuencia de corte en $\text{Hz}$ del filtro pasa-bajos que precede al diezmado por $M=3$?
    \end{itemize}
\end{enumerate}

\section{Actividad 3: Remuestreo Racional y Resolución ($\mathbf{L/M}$)}

Configure la aplicación en modo \textbf{Remuestreo (L/M)} con los siguientes parámetros:
\begin{itemize}
    \item $F_s = 200 \text{ Hz}$
    \item $f_0 = 20 \text{ Hz}$
    \item Factor $L = 5$
    \item Factor $M = 4$
\end{itemize}

\subsection{Preguntas y Tareas}

\begin{enumerate}
    \item \textbf{Nueva Tasa:} Calcule el nuevo factor de remuestreo $R = L/M$. ¿Cuál es la nueva frecuencia de muestreo $F_s'$?
    \item \textbf{Filtro Único:} El proceso de remuestreo $L/M$ utiliza un único filtro con $\omega_c = \min\{\pi/L, \pi/M\}$.
    \begin{itemize}
        \item ¿Cuál de los dos factores ($L$ o $M$) domina la elección de la frecuencia de corte en este caso? ¿Cuál es la frecuencia de corte en $\text{Hz}$?
    \end{itemize}
    \item \textbf{Zero-Padding (Resolución DFT):} Desactive y reactive la casilla \textbf{Aplicar Zero-Padding a la DFT} en la salida (AX4).
    \begin{itemize}
        \item Al desactivar el ZP, ¿se mueve el pico de $20 \text{ Hz}$? ¿Qué cambia en el espectro (AX4)?
        \item Si la señal original de entrada $x[n]$ fue muestreada por $N=128$ puntos, ¿cuál es la resolución de frecuencia $\Delta f$ real (la capacidad de distinguir tonos) de la DTFT? (Depende solo de $F_s$ y $N$).
    \end{itemize}
\end{enumerate}

\end{document}